\documentclass[a4paper,10pt]{book}

% 引入我们自定义的样式包
\usepackage{dag-notes}
\usepackage{fontawesome5}
\usepackage{CJKutf8}
\usepackage{tikz-cd}
\usepackage{geometry}
\geometry{top=1.5cm, bottom=1.5cm, left=1.5cm, right=1.5cm} % 设置适合 12 页阅读的页边距



\newcommand{\OG}{\operatorname{\mathcal{O}_G}}
\newcommand{\Top}{\operatorname{Top}}
\newcommand{\Hom}{\operatorname{Hom}}
\newcommand{\inHom}{\underline{\operatorname{Hom}}}
\newcommand{\Ab}{\operatorname{Ab}}
\newcommand{\C}{\mathscr{C}}
\newcommand{\D}{\mathcal{D}}
\newcommand{\op}{\mathrm{op}}
\newcommand{\catinf}{\mathrm{Cat}_\infty}
\newcommand{\infcat}{\infty\text{-category}}
\newcommand{\infcats}{\infty\text{-categories}}
\newcommand{\CAlg}{\operatorname{CAlg}}
\newcommand{\Sch}{\mathrm{Sch}}
\newcommand{\Etale}{\text{Étale}}
\newcommand{\QCoh}{\mathrm{QCoh}}
\newcommand{\qc}{\mathrm{qc}}
\newcommand{\qcqs}{\mathrm{qcqs}}


\title{Six-Functor Formalisms over Orbit Categories}
\author{
}
\date{\today}
\begin{document}
% 生成标题
\maketitle
% 导入各个章节
% 使用 \input 可以保持紧凑排版（不会强制换页）
% 如果想要每一章从新的一页开始，可以使用 \include


\tableofcontents

\chapter{Introduction}

Our Main idea come from \cite{scholzeSFF} and \cite{scholzeGHT}. 







\input{sections/Stable Homotopy Theory}

\chapter{Equivariant Homotopy Theory and Orbit Category}

Let $G$ denote a compact Lie group. 

\begin{definition}
    The orbit category $\OG$ associated to $G$ is defined by 
    \begin{enumerate}
        \item[Objects] $\{G/H\mid H\text{ is a closed subgroup of }G\}$
        \item[Morphsims] Continuous $G$-equivariant maps. 
    \end{enumerate}
\end{definition}

Let $\Hom_{\OG}(G/H,G/K)$ equipped with compact-open topology, then we have the homeomorphsim 
\begin{align*}
    \Hom_{\OG}(G/H,G/K) \simeq (G/K)^H.
\end{align*} 




\begin{definition}
    An $\OG$-space is a contravariant functor from $\OG$ to $\Top$.
\end{definition}

\begin{definition}
    A $G$-space is a covariant functor from classifying category $BG$ to $\Top$.
\end{definition}

\begin{theorem}[Elmendorf's Theorem]
    
\end{theorem}






\chapter{Six-Functor Formalisms}

Throughout this chapter, we let $\C$ denote a category satisfying the following hypotheses
\begin{itemize}
    \item[a.] The category $\C$ has all finite limits 
    \item[b.] The class of morphisms $E$ contains all isomorphisms. $E$ is stable under pullback, composition and diagonals.
\end{itemize} 

A six-functor formalism is a functor $\D:\C^{\op} \to \CAlg(\catinf)$ sending any morphism $f$ in $
\C$ to a symmetric monoidal functor $f^*$ in $\CAlg(\catinf)$ and an assignment
\begin{align*}
    f:X\to Y \qquad \mapsto \qquad f_!:\D(X)\to \D(Y)
\end{align*}
for any $f\in E$ to a functor of symmetric monoidal $\infcats$.\footnote{$f_!$ need not be a symmetric monoidal functor.}
We call functor $f^*$ the star pullback and $f_!$ shriek pushforward.

Additionally, we require that 
\begin{itemize}
    \item[1] There exists functor $f_*$ right adjont to $f^*$, we call $f_*$ star pushforward;
    \item[2] There exists functor $f^!$ right adjont to $f_!$, we call $f^!$ shriek pullback;
    \item[3] For any $X\in\C$ and $A\in \D(X)$, the symmetric monoidal $\infcat$ $\D(X)$ admits a adjoint pair of functor 
    \begin{align*}
         - \otimes A \dashv \inHom(A,-),
    \end{align*} 
    we call it tensor product and internal Hom;
    \item[4] The star pullback functor $f^*:\D(Y)\to \D(X)$ compatible with symmetric monoidal structure on $\D(X)$ and $\D(Y)$ and compatible with composition;
    \item[5] The shriek pushforward functor $f_!: \D(X)\to \D(Y)$ compatible with composition and satisfying the base change and projection formula isomorphism.
\end{itemize}
 
Here the base change formula means if there exists a Cartesian square 
\[
\begin{tikzcd}
X' \arrow[r, "g'"] \arrow[d, "f'"'] 
  \arrow[dr, phantom, "\lrcorner", very near start]
& X \arrow[d, "f"] \\
Y' \arrow[r, "g"'] & Y
\end{tikzcd}
\]
over $\C$, then the diagram 
\[
\begin{tikzcd}
D(X') \arrow[d, "f'_!"'] 
& D(X) \arrow[l, "g'^*"'] \arrow[d, "f_!"] \\
D(Y') 
& D(Y) \arrow[l, "g^*"]
\end{tikzcd}
\]
commutes up to homotopic coherent or equivariantly we say $f_!'\circ g'^*$ is natural isomorphic to $g^*\circ f_!$.

The projection formula means $\D(X)$ is equipped with $\D(Y)$-module structure via $f^*:\D(Y)\to \D(X)$ and functor $f_!$ is $\D(Y)$-linear. Equivalently, for any $A\in X$ and $B\in Y$ the shriek pushforward satisfies
\begin{align*}
    f_!(A \otimes f^*B) = f_!(A)\otimes B.
\end{align*}

\begin{example}[Six functor formalism for abelian sheaves]
    Let $S$ be the base scheme and $\tau$ denote a Grothedieck topology of $\Sch_{/S}$.\footnote{We can choose $\tau$ be the Zariski, \Etale, Nisnevich, fpqc, or fppf topology.}
    Define the functor $\D:(\Sch_{/S})^\op \to \CAlg(\catinf)$ by assigning
    \begin{align*}
        X \mapsto \D(\Ab(X_{\tau})) \qquad (f:X\to Y) \mapsto (f^{-1}:\D(\Ab(Y_{\tau}))\mapsto \D(\Ab(X_{\tau})))
    \end{align*}
    where $X\in \Sch_{/S}$ and $\D(\Ab(X_{\tau}))$ denotes the derived $\infcat$ of  the abelian sheaves on $X$.
    Let $E$ be the class of compactifiable morphisms, i.e., for any $f:X\to Y$ in $E$, there exists a factorization
    \begin{align*}
        f:X\xrightarrow{j} \bar{X} \xrightarrow{p} Y
    \end{align*} 
    such that $j$ is an open immersion and $p$ is a proper morphism. Define $f_!=  Rp_*\circ j_!$ where $j_!$ is extension by zero.  
\end{example}

\begin{example}[Six functor formalism for quasi-coherent sheaves]
    Let $S$ be the base scheme. Define the functor $\D:(\Sch_{/S})^\op \to \CAlg(\catinf)$ by assigning
    \begin{align*}
        X\mapsto \D_{\qc}(X) \qquad f \mapsto Lf^*.
    \end{align*}
    where $\D_{\qc}(X)$ denotes the derived $\infcat$ of complexes with quasi-coherent cohomology sheaves on $X$.\footnote{When $X$ is qcqs, we have $\D_{\qc}(X)=\D(\QCoh(X))$.}
    Let 
    \begin{align*}
        E =\{\text{proper morphism in }(\Sch_{/S})_{\qcqs}\}
    \end{align*}
    Define the shriek pushforward by $f_!:= Rf_*$. The existence of the right adjoint of $f_!$ is the statement of Grothedieck duality.
\end{example}











\input{sections/Gestalt}

\bibliographystyle{alpha}
\bibliography{ref}
\end{document}